% ============================================================================
% Dutch UI and structural strings for the One Physics Book series.
% Loaded by styles/onephysics.sty when \booklang is nl.
% ============================================================================

% Theorem / statement environment titles
\newcommand{\omnameDefinition}{Definitie}
\newcommand{\omnameTheorem}{Stelling}
\newcommand{\omnameProposition}{Propositie}
\newcommand{\omnameLemma}{Lemma}
\newcommand{\omnameCorollary}{Gevolg}
\newcommand{\omnameMethod}{Methode}
\newcommand{\omnameExample}{Voorbeeld}
\newcommand{\omnameNotation}{Notatie}
\newcommand{\omnameRemark}{Opmerking}
\newcommand{\omnameExercise}{Oefening}
\newcommand{\omnameExercises}{Oefeningen}
\newcommand{\omnameProblem}{Probleem}
\newcommand{\omnameProblems}{Problemen}
\newcommand{\omnameProof}{Bewijs}

% Admitted results and solutions appendix headers
\newcommand{\omadmittedtext}{Toegegeven op dit niveau.}
\newcommand{\omsolutionof}{Oplossing van}
\newcommand{\ompage}{pagina}
\newcommand{\omnameGotoSolution}{Oplossing p.}

% Misc text macros
\newcommand{\st}{\text{ zodat }}

% Cover / title page
\newcommand{\bookmaintitle}{One Physics Book}
\newcommand{\booktagline}{Het natuurkundeboek om ze allen te regeren}
\newcommand{\bookdescription}{Natuurkunde van de kleuterschool tot het einde van de bachelor}
\newcommand{\bookcontributors}{De bijdragers van het One Physics Book}
\newcommand{\bookversionlabel}{Versie}

% Colophon
\newcommand{\omcolophonsource}{Broncode, issues en bijdragen:}

% Solutions appendix part title
\newcommand{\omsolutionspart}{Oplossingen van de oefeningen}

% Part titles (Jaren 1–12; parallel to English Grade N)
\@namedef{omstr@part.grade1}{Jaar 1 (6--7 jaar)}
\@namedef{omstr@part.grade2}{Jaar 2 (7--8 jaar)}
\@namedef{omstr@part.grade3}{Jaar 3 (8--9 jaar)}
\@namedef{omstr@part.grade4}{Jaar 4 (9--10 jaar)}
\@namedef{omstr@part.grade5}{Jaar 5 (10--11 jaar)}
\@namedef{omstr@part.grade6}{Jaar 6 (11--12 jaar)}
\@namedef{omstr@part.grade7}{Jaar 7 (12--13 jaar)}
\@namedef{omstr@part.grade8}{Jaar 8 (13--14 jaar)}
\@namedef{omstr@part.grade9}{Jaar 9 (14--15 jaar)}
\@namedef{omstr@part.grade10}{Jaar 10 (15--16 jaar)}
\@namedef{omstr@part.grade11}{Jaar 11 (16--17 jaar)}
\@namedef{omstr@part.grade12}{Jaar 12 (17--18 jaar)}

% Part titles (Bachelorjaren 1–3; parallel to English Bachelor Year N)
\@namedef{omstr@part.bachelor1}{Bachelor jaar 1 (18--19 jaar)}
\@namedef{omstr@part.bachelor2}{Bachelor jaar 2 (19--20 jaar)}
\@namedef{omstr@part.bachelor3}{Bachelor jaar 3 (20--21 jaar)}

% Plural names + \omnameChapter, used by the \crefname declarations in
% onephysics.sty. Without these, cleveref falls back to its English built-ins.
\newcommand{\omnameDefinitions}{Definities}
\newcommand{\omnameTheorems}{Stellingen}
\newcommand{\omnamePropositions}{Proposities}
\newcommand{\omnameLemmas}{Lemma's}
\newcommand{\omnameCorollarys}{Gevolgen}
\newcommand{\omnameMethods}{Methoden}
\newcommand{\omnameExamples}{Voorbeelden}
\newcommand{\omnameNotations}{Notaties}
\newcommand{\omnameRemarks}{Opmerkingen}
\newcommand{\omnameChapter}{Hoofdstuk}
\newcommand{\omnameChapters}{Hoofdstukken}

% LaTeX's own structural names. babel would normally localise these, but it is
% optional here (see onephysics.sty).
\renewcommand{\chaptername}{Hoofdstuk}
\renewcommand{\partname}{Deel}
\renewcommand{\contentsname}{Inhoudsopgave}
\renewcommand{\appendixname}{Bijlage}
\renewcommand{\indexname}{Index}
